\documentclass[]{../../../../tex/classes/styledArticle}
\usepackage{../../../../tex/packages/hebrewSupport}
\usepackage{../../../../tex/packages/mathShortcuts}
\usepackage{../../../../tex/packages/theoremsSupport}

\usepackage{listings}
\lstset{
	language=Python,
	numbers=left,
	breaklines,
	basicstyle=\ttfamily,
	commentstyle=\color{gray},
	prebreak=\textrightarrow,
}


\author{שחר פרץ}
\title{אלגוריתמים $\sim$ \textit{הרצאה }}
\begin{document}
	\maketitle
	\textbf{מרצה: }
	
	\subsection{מרחקים קצרים ביותר בין כל הזוגות}
	איך נמצא מרחקים קצרים ביותר בין זוג צמתים? 
	\begin{itemize}
		\item דייקסטרא מכל קודקוד, ורק עם משקלים אי־שליליים. סיבוכיות $\oc(\sof V^2 \log V + \sof V \sof E)$
		\item אם יש משקלים שליליים, נוכל להריץ BF מכל קודקוד ולקבל פתרון ב־$\oc(\sof V^2 \sof E)$. 
	\end{itemize}
	
	איך בכלל נרצה לפלוט דברים? באמצעות מטריצה. מטריצה אחת תחזיק את המרחקים בין הקודקוד ה־$i$ לקודקוד ה־$j$ (נסמן $D_{ij}$), ו־$P_{ij}$ שמחזיר צומת אחד לפני במק''ב בין $v_i$ ל־$v_j$. נסמן את הקודקודים להיות $V = \{v_1 \dots v_n\}$. ואז המטריצות $D, P \in M_{n \times n}$. 
	
	מעצם העובדה שהפלט בגודל $\sof V^{2}$, אין מניעה (מנקודת מבט של סיבוכיות זכרון) שלא נקבל/נעבד את הגרף להיות מטריצת שכנויות $W$, עם שינוי קל: 
	\[ W_{ij} = \begin{cases}
		0 & i = j \\
		w(i, j) & (i, j) \in E \\
		\infty & \other
	\end{cases} \]
	נרצה בדומה ל־BF למצוא לאט לאט את המסלול מאורך $n$ הקל ביותר (נסמן את המטריצה בשלב ה־$i$ ב־$D^{(i)}$).  
	נבחין שעבור $n = 1$, $D^{(1)} = W$ (המסלול באורך $1$ הקצר ביותר ניתן ע''י המטריצה $W$). 
	מכאן שאם נדע לחשב את $D_{ij}$ בהינתן $D_{ij}^{(n - 1)}$, פרקטית סיימנו. נבחין ש־: 
	\[ D^{(m + 1)}_{ij} = \min_{k}\{D_{ik}^{(m)} + W_{kj}\} \]
	
	לגבי זיהוי מעגלים שליליים: לאחר $n$ איטרציות, נוכל: 
	\begin{itemize}
		\item להשתמש בדרך הנאיבית של BF, שזה לעשות עוד איטרציה ולבדוק האם המטריצה משתנה. 
		\item דרך יעילה יותר, היא לבדוק (אפשר בכל איטרציה ואפשר רק באחרונה, זניח מבחינת סיבוכיות ויכול לעצור את האלגוריתם מוקדם) האם $D_{ii}^{n} < 0$, ובמקרה זה נסיק שיש מעגל שלילי. 
	\end{itemize}
	יש לזה רק חסרון אחד, והוא $\oc(\sof V^{4})$. 
	
	נבחין שעדכון $D^{(m + 1)}$ לפי $D^{(m)}$ ו־$W$ ממש מזכיר כפל מטריצות (עברו על הנוסחה לעיל), עם השינויים הבאים: חיבור במקום כפל, ומינימום במקום חיבור (סכום). ב''כפל'' זה, $D^{(n)} = W^{(n)}$. חזקה (של כל פעולה) אפשר לעשות ב־$\log n$ כפול העלות של כפל (במקרה הזה – $\sof V^3 \log \sof V$), באמצעות interated squaring. 
	\rmark{לא צריך לעבור לייצוג בינארי. אין שום דבר שמונע מאיתנו לחשב את $D^{2^{k}}$ בעבור $2^{k} > n$ במקום את $D^{n}$. זה מפשט את האלגוריתם. }
	
	ננסה לשפר עוד. נגיע לאלגוריתם של Floyd-Warshal: במקום לעדכן מסלולים לפי כמות הקשתות, באיטרציה ה־$k$ נרשה רק מסלולים מגובה $k$, כלומר כאלו שהצמתים שלהם משתתפים רק $v_1, v_2 \dots v_k$. ועכשיו, באיטרציה ה־$k$, לא נעשה חיפוש על המינימום של כל ה־$k$־ים, אלא רק נשווה את המסלול הכי קצר מגובה עד $k$ (שחושב באיטרציה הקודמת) או מ־$k + 1$ (חישוב $\oc(1)$). זה מוריד את הסיבוכיות של כל פעולה לכדי $\oc(\sof V^{2})$, ומאפשר לנו להפעיל את האלגוריתם ב־$\oc(\sof V^{3})$. לא נוכל להשתמש ב־iterated squring במקרה זה. באלגוריתם הזה, העדכון והבסיס יפעלו כך: 
	\begin{align*}
		\dg_0(i, j) = W_{ij} = \begin{cases}
			w(i, j) & (i, j) \in E \\
			0 & i = j \\
			\infty & \other
		\end{cases} && \dg_{k + 1}(i, j) = \min\{\dg_k(i, j), \dg_k(i, k + 1) + \dg_k(k + 1, j)\}
	\end{align*}
	
	נבחין שהשגנו שיפור אסימפטוטי (במיוחד בגרפים צפופים) על פני BF בעבור אלגוריתם שהוא אשכרה אלגנטי ומהיר למימוש. ההוכחה שהנוסחה לעיל נכונה ושומרת על השמורה (הגדרה) של $D$, היא: 
	\begin{proof}
		המלסול הכי טוב בין $v_i$ ל־$v_j$ מגובה $k + 1$ הוא: 
		\begin{itemize}
			\item או מסלול מגובה לכל היותר $k$, ואז כבר מצאנו אותו. 
			\item או מסלול מגובה $k + 1$, כלומר שרשור של המסלול מגובה $k$ בין $v_i$ ל־$v_{k + 1}$ ומסלול מגובה בדיוק $k$ בין $v_{k + 1}$ ל־$v_j$. 
		\end{itemize}\envendproof
	\end{proof}
	\rmark{מסלול מגובה $k$ אומר לכל היותר $k$. }
	\subsection*{Johnson Division}
	השמועה אומרת שאוניקס טרוניקס הפסידו בבית הזה. 
	
	האלגוריתם של ג'ונסון מנסה לעשות ''דייקסטרא מכל קודקוד``. על פניו, צריך משקלים אי־שליליים. בעפ''מים, לא הייתה לנו בעיה להוסיף קבוע לכל המשקלים, כי כמות הקשתות הייתה קבועה. אבחנה טובה היא שאחרי שהעלנו את המשקלים, אכפת לנו רק על זה שלא שינינו את הסדר בין מסלולים בין אותם הקודקודים. 
	
	\theo{תהא $h \co V \to \R$ פונקציה כלשהי. נגדיר פונקציית משקל חדשה $w^*(u, v) = w(u, v) + h(u) - h(v)$. אזי, מסלול $v \squig u$ הוא מק''ב לפי $w$ אמ''מ הוא מק''ב לפי $w^*$. }\begin{proof}
		עבור מסלול $u = v_0 \to v_1 \to v_2 \to \cdots \to v_k = v$, מתקיים: 
		\[ \mathrm{cost}_w(u, v) = \sum_{i = 1}^{k} w(v_{i - 1}, v_i)   \]
		בעוד: 
		\[ \mathrm{cost}_{w^*}(u, v) = \sum_{i = 1}^{k}w(v_{i - 1}, v_j) + h(v_{i - 1}) + h(v_i)) = \sum_{i = 1}^{k} w(v_{i - 1}, v_i) + \sum_{i = 1}^{k} + h(u) - h(v) \]
		כלומר, כל המק''בים $v \squig u$ שינו את המשקל שלהם בדיוק ב־$h(u) - h(v)$ ולכן הטענה נכונה. 
	\end{proof}
	
	עתה, רק נותר למצוא $h$ כזו כך שהיא הופכת את משקלי הקודקודים לאי־שליליים (כלומר $\forall e \in E \co w^*(e) \ge 0$), כדי שנוכל להפעיל דייקסטרא ולגמור עניין. 
	אז מה בעצם אנו רוצים לעשות: 
	\[ \forall (u, v) \in E \co w(u, v) + h(u) - h(v) = w^*(u, v) \ge 0 \]
	הביטוי לעיל מתקיים אמ''מ: 
	\[ \forall (u, v) \in E \co h(v) \le h(u) + w(u, v) \]
	הא''ש לעיל נראה בדיוק כמו א''ש המשולש, מסיבה כלשהי. ננסה למצוא פונקציית ''מרחק``. 
	
	ניעזר בטריק האהוב שלנו מההרצאה הקודמת קודמת. נוסיף קודקוד חדש עם קשתות מכוונות (ממשקל $0$) אל כל אחד מהקודקודים האחרים ונקווה לטוב: נחשב בעזרת BG את $\dg(s, v)$ לכל צומת. נגדיר $h(v) = \dg(s, v)$. עכשיו $w^*(u, v) \ge 0$ ונריץ דייקסטרא מכל קודקוד. 
	
	לגבי הסיבוכיות: 
	\begin{itemize}
		\item הוספת צומת והרצת BF: $\oc(\sof V \sof E)$. 
		\item דייקסטרא מכל צומת: העלות הרגילה של $\oc(\sof V^2 \log \sof V  + \sof V \sof E)$. 
	\end{itemize}
	
	נעבור קצת לשאלות תרגול. 
	
	\section*{לא תרגול}
	\defi{\textit{סגור טרנזטיבי} של גרף זהו גרף עם אותם קודקודים בדיוק, ויש קשת $(u, v)$ בגף אמ''מ בגרף המקורי יש מסלול $u \squig v$. }
	\exe{איך מוצאים סגור טרנזטיבי של גרף? }
	הגישה הבנאלית – לחשב את המרחקים בין כל הזוגות ולשים קשת אם המרחק לא אינסוף. הגישה העוד יותר בנאלית – לעשות BFS מכל אחד (זה יעלה פחות). זו גם הגישה הכי טובה. 
	
	אולי מזה אנחנו לומדים שלפעמים אנו מקובעים על אלגוריתם אחד, אבל משהו פשוט יותר (ובלי פיבונאצי באמצע) דווקא הוא זה שיגרום לדברים לעבוד. 
	
	\exe{רוצים לדעת אם יש מסלול (לאו דווקא פשוט) באורך בדיוק $m$ בין כל זוג קודקודים (אין משקלים). }
	ניעזר בווריאציה על האלגוריתם שלנו מתחילת ההרצאה: נגדיר $D_{ij}^{(m)}$ להיות ''קיים מסלול באורך בדיוק $m$ בין $i$ ל־$j$``. לגבי הבנייה הרקורסיבית, הבסיס יראה כך: 
	\begin{align*}
		D_{ij}^{(1)} = \begin{cases}
			1 & (i, j) \in E \\
			0 & \other 
		\end{cases} && D_{ij}^{(m + 1)} = \bigvee_k\cl{D_{ij}^{(m)} \land D_{kg}^{(1)}}
	\end{align*}
	כאשר המשמעות של ה''$\land, \lor$`` היא and במקום כפל ו־or במקום חיבור. נוכל לשפר את הסיבוכיות עם iterated sqauring, ולקבל סה''כ $\oc(\sof V^{3}\log m)$. 
	יש כאן משהו יפה שנוכל לעשות: אם נעביר את ה־and להיות חזרה כפל, ו־or להיות חזרה חיבור, נקבל ב־$D_{ij}$ את כמות המסלולים (הלא פשוטים) באורך בדיוק $m$ (ו־$0$ אם אין). זה גם משפר את הסיבוכיות הודות לאלגוריתמים יעילים לכפל מטריצות. 
	
	\exe{נתון גרף מכוון $G = (V, E)$, ומשקל $w \co E \to \R$. נתון מספר $k$. נרצה למצוא את כל הקשתות שנמצאות על מעגל פשוט עם משקל קטן או שווה ל־$k$. }
	הפתרון מאוד פשוט. נמצא את גרף המ''קבים (מסלולים קלים ביותר), ועכשיו לכל קשת $e = (u, v)$ נבדוק האם המק''ב $v \squig u$ הוא ממשקל לכל היותר $k - w(u, v)$ (שימו לב שזה עובד רק אם הגרף מכוון). 
	
	
	\textbf{חידה: }יש בניין עם 100 קומות ושתי ביצים קשות (ממש קשות). רוצים לגלות מאיזה קומה הן מתחילות להשבר. יש להן פיצ'ר לפיו אם ביצה נשברת, הביצה נשברה ואי אפשר לשבור אותה שוב. יש לך 2 ביצים. \textbf{השאלה: }אסטרטגייה לזריקת ביצים שהיא הכי טובה ב־WC ביחס למספר זריקות. 
	
	
	
	
	
	\ndoc
\end{document}